\documentclass[a5paper,10pt]{article}

% https://github.com/InDevRus/filippov-solutions

% Нумерация страниц
\pagenumbering{gobble}

% Настройка полей
\usepackage{geometry}
\geometry{left=1cm}
\geometry{right=1cm}
\geometry{top=1cm}
\geometry{bottom=1.5cm}

% Вставка таблицы
\usepackage{array}
\usepackage{tabularx}

% Инструменты выравнивания формул
\usepackage{amsmath}
\usepackage{mathtools}

% Диакритика в математике
\usepackage{amsfonts}

% Вычисления размеров на ходу
\usepackage{calc}

% Удобные записи производных
\usepackage[italic=true]{derivative}

% Настройки сносок
\usepackage{enotez}

% Длинные знаки векторов
\usepackage{anyfontsize}
\usepackage[b]{esvect}

% Вставка картинок
\usepackage{graphicx}

% Вставка красной строки
\usepackage{indentfirst}
\setlength{\parindent}{1em}

% Условия в командах
\usepackage{xifthen}

% Луа-скрипты
\usepackage{luacode}

% Настройка команд отдельно в математическом и текстовом режимах
\usepackage{mathcommand}

% Для повторения LaTeX кода
\usepackage{multido}

% Конфигурация отступов формул
\delimitershortfall-1sp
\usepackage{mleftright}
\mleftright

% Растягивание объектов
\usepackage{relsize}

% Обтекание текстом
\usepackage{wrapfig}
\usepackage{subcaption}
\usepackage{floatflt}

% Поддержка ввода кириллицы
\usepackage[english, russian]{babel}

% Настройка корневой позиции для компилятора
\usepackage{import}
\providecommand*{\ProjectRootPath}{..}

\import{\ProjectRootPath/preambles/macros/}{theorems}

\import{\ProjectRootPath/preambles/fonts}{font_setup}

% Знак градуса
\usepackage{siunitx}

\import{\ProjectRootPath/preambles/macros/}{operators}
\import{\ProjectRootPath/preambles/macros/}{redefinitions}

\import{\ProjectRootPath/preambles/fonts}{letters_setup}
\import{\ProjectRootPath/preambles/fonts/adjustments}{custom_superscripts}

\import{\ProjectRootPath/preambles/custom}{formatting}
\import{\ProjectRootPath/preambles/custom}{frames}
\import{\ProjectRootPath/preambles/custom}{list_setup}

% TODO: перевести команды на современный стиль объявления

\begin{document}
    $\br{x - z} \parder {z} {x} + \br{y - z} \parder {z} {y} = 2 z$.
    $\dfrac {dx} {x - z} = \dfrac {dy} {y - z} = \dfrac {dz} {2z}$.
    Интегрируемыми комбинациями являются
    $\dfrac {dx} {x - z} = \dfrac {dz} {2z}$
    и
    $\dfrac {dy} {y - z} = \dfrac {dz} {2z}$. Они равносильны, так что достаточно решить только одну. К ним также можно добавить функцию $z\br{x; y} = 0$, хотя это решение и не подходит под условие задачи. \vspace{1mm}

    Решим, например, первую.
    $\dfrac {dx} {dz} = \dfrac {1} {2z} x\br{z} - \dfrac {1} {2}$.
    Решение этого линейного уравнения первого порядка нужно рассматривать в областях $z > 0$ и $z < 0$ по отдельности.

    \begin{itemize}
        \item Рассмотрим уравнение выше в области $z > 0$.

        $\dfrac {1} {\sqrt{z}} \dfrac {dx} {dz} - \dfrac {1} {2\sqrt{\pow{z}{3}}} x = -\dfrac {1} {2 \sqrt{z}}$.
        $\dfrac {d} {dz} \br{\dfrac {1} {\sqrt{z}} \cdot x} = -\dfrac {1} {2 \sqrt{z}}$.
        $\dfrac {x} {\sqrt{z}} = -\sqrt{z} + A_{1}$.
        $\dfrac {x + z} {\sqrt{z}} = A_{1}$.

        \item В области $z < 0$ замена переменных $- z = \zeta > 0$, $\dfrac {dx} {dz} = -\dfrac {dx} {d\zeta}$
        приводит к уравнению
        $\dfrac {dx} {d\zeta} - \dfrac {1} {2\zeta} x\br{\zeta} = \dfrac {1} {2}$,
        откуда $\dfrac {x - \zeta} {\sqrt{\zeta}} = A_{2}$,
        то есть $\dfrac {x + z} {\sqrt{-z}} = A_{2}$.
    \end{itemize}

    Таким образом, при $z > 0$ имеем первые интегралы
    $\dfrac {x + z} {\sqrt{z}} = A_{1}$ и $\dfrac {y + z} {\sqrt{z}} = B_{1}$;
    при $z < 0$ имеем $\dfrac {x + z} {\sqrt{-z}} = A_{2}$ и $\dfrac {y + z} {\sqrt{-z}} = B_{2}$.

    Переходим к частной поверхности. Её параметризация по параметру $t$ имеет вид $x\br{t} = t$, $y\br{t} = t - 2$, $z\br{t} = 1 - 2t$.

    Сначала рассмотрим область пространства $z > 0$, $t < \dfrac {1} {2}$.
    Подставив параметризацию, получим
    $\dfrac {1 - t} {\sqrt{1 - 2t}} = A$, $\dfrac {-1 - t} {\sqrt{1 - 2t}} = B$.
    С одной стороны, имеем  $\dfrac {1 - t} {-1 - t} = \dfrac {A} {B}$, а с другой $\dfrac {2} {\sqrt{1 - 2t}} = A - B$.
    Выразив параметр $t$, исключим его следующим образом:
    $t = \dfrac {-\br{A + B}} {A  - B} = \dfrac {\br{A - B}^2 - 4} {2\br{A - B}^2}$.
    Упростив последнее, получим $\br{3A + B}\br{A - B} = 4$.
    Итак, искомая поверхность задаётся равенством $\br{3x + y + 4z}\br{x - y} = 4z$.

    В области $z < 0$, $t > \dfrac {1} {2}$ получаем
    $t = \dfrac {-\br{A + B}} {A  - B} = \dfrac {\br{A - B}^2 + 4} {2\br{A - B}^2}$,
    а соответственно и $\br{3A + B}\br{A - B} = -4$.
    Окончательно, $\br{3x + y + 4z}\br{x - y} = 4z$.

    Поскольку тождества одинаковы, поверхность
    $\br{3x + y + 4z}\br{x - y} = 4z$ при $z \ne 0$ является ответом.

\end{document}
